\documentclass[11pt,reqno]{amsart}
\usepackage[T1]{fontenc}
\usepackage{amsmath,amssymb,amsthm,mathtools}
\usepackage{geometry,booktabs}
\usepackage[colorlinks=true,linkcolor=blue,citecolor=blue,urlcolor=blue]{hyperref}
\geometry{margin=1in}
\numberwithin{equation}{section}
\newtheorem{theorem}{Theorem}[section]
\newtheorem{lemma}[theorem]{Lemma}
\newtheorem{proposition}[theorem]{Proposition}
\newtheorem{corollary}[theorem]{Corollary}
\theoremstyle{definition}
\newtheorem{definition}[theorem]{Definition}
\newtheorem{remark}[theorem]{Remark}
\newtheorem{question}[theorem]{Open question}
\newcommand{\C}{\mathbb C}
\newcommand{\R}{\mathbb R}
\newcommand{\Q}{\mathbb Q}
\newcommand{\W}{\mathcal W}
\newcommand{\T}{\mathcal T}
\newcommand{\Om}{\Omega}
\newcommand{\eps}{\varepsilon}
\newcommand{\pp}{\mathrm{pair}}
\DeclareMathOperator{\comp}{component}
\title[Real primitive tips: conditional geometry]{Real primitive-tip returns and conditional marked-Wave bounds:\\ a referee-controlled revision}
\author{Research compilation from loops-research-archive}
\date{16 November 2026 (AI-assisted; not independently reviewed)}
\begin{document}
\begin{abstract}
We revisit a selected real near-$-2$ family of primitive-tip return
candidates, with exact finite root packets and a written uniform
$p\ge13$ tail argument for local parameter isolation, orbit order and
degree-two polynomial-like returns.  No global real-root uniqueness
is asserted.  The original manuscript further announces an exact
KKL physical marking, a marked-width sandwich and
$D_p=2/C_p-1/P_p>1/100$ for every $p\ge3$.
The exact verifiers check rational inequalities and finite
incidence predicates, \emph{not} the continuum identification,
primeness transport, full KKL A.1 proper-path classes or
Schwarz--Christoffel admissibility.  We therefore retain the
strict numerical consequences only \emph{conditionally on the
explicit geometric bridges (H1)--(H3)}.  Likewise, exclusion of
outer-incident buffered canonical WAD weights for $p\ge11$
is conditional on the cyclic-cover interpretation (H4).
The six remaining inner endpoint types, a strictly positive
buffered coefficient, a four-push limit and uniform
marked quasiconformal depth transport are not settled.
We give finite depth witnesses but no unconditional
suffix-straightening claim.  The accompanying hostile
referee report records source-specific gaps and the
verification boundary.  This AI-assisted archive revision
is not independent mathematical peer review and
does not claim new Mandelbrot-set local connectivity.
\end{abstract}
\maketitle
\tableofcontents
\section{Introduction and proof status}
This paper studies a \emph{real}, selected, postcritically finite quadratic family.  The period-three tip belongs to the $p$-return, not to the original polynomial of period $p$.  In particular $x_{jp}=f^{jp}(0)$, not $x_j=f^j(0)$.  The results concern the physical tree slit and its central degree-two pullback, not an arbitrary puzzle or a global Mandelbrot statement.  KKL \cite{KKL2026} is a posted research preprint, not a claimed peer-reviewed journal publication.

The stated unconditional algebraic and local-return arguments are distinguished from the \emph{conditional} global marked-geometry deductions.  Earlier claimed proofs of the missing bridges are preserved for checking, not endorsed by a passing finite verifier.  Each derivation distinguishes exact rational interval computations, written complex-analytic deductions, and named external theorems.  Passing a Python or Lean check certifies the finite premises it implements but never formalizes KKL, straightening, Carath\'eodory boundaries, or conformal extremal width.  The complete verification commands and precision are in \S\ref{sec:reproduce}.  The copied \texttt{REVIEW\_FINDINGS.md} records the original audit's assertions; the independent root-level \texttt{REFEREE\_REPORT.md} supersedes its assessment wherever they differ.  The mathematical authorship and independent review of this synthesis remain to be established.

\section{Referee-controlled status and conditional geometric inputs}\label{sec:status}

The term \emph{proved} below is reserved for consequences for which the
stated input hypotheses and the displayed argument suffice.  An exact
rational verifier checks a finite predicate; it does not certify a
topological or analytic identification.  The companion independent
\texttt{REFEREE\_REPORT.md} (at the repository root) explains why the
following bridges must not be suppressed in an unconditional statement.
The substantial calculations of the original manuscript are preserved
in later sections as \emph{conditional derivations and proof attempts};
their mathematical existence in the archive is not independent review.

\begin{definition}[Unresolved geometric bridges]\label{def:bridges}
For a selected period $p$, we distinguish the following hypotheses
from the exact polynomial and rational inequalities.
\begin{description}
\item[(H1) Renormalization/nesting.]
The central $3p$-return has the asserted Chebyshev filled interval
with the two named endpoints, and any hypothetical competing
period-two renormalization of the \emph{central quadratic-like map}
is nested so that McMullen's almost-disjointness theorem can be
transported to its two little Julia sets; both inverse critical
points in the primeness contradiction belong to the proper
degree-two return domain.  No regularity of $\partial V_p$
stronger than local connectedness is silently assumed.

\item[(H2) Physical KKL ideal marking.]
On the specific two slit complements, the dynamically defined
KKL Hubbard tree, the hyperbolic prime-end/ideal boundary
marking, and the complete proper-path and subordination relation
are identified with the stated physically indexed slit-side
accesses.  In particular the inclusion's pulled marking
is compatible with the source covering marking and with
the refinement required by KKL Proposition~A.12.
This is a \emph{continuum class-identification} condition,
not merely a ten/twenty integer table.

\item[(H3) Full-path conformal comparison.]
The proposed embedded parent Schwarz--Christoffel
quadrilateral foliations consist of admissible simple
KKL A.1 Waves with their \emph{minimal} prescribed
bases and crests.  Every admissible source child Wave
has a lift into the proposed entire four-electrode
all-path envelope; the square-sum periodization is
admissible for \emph{all} such leaves, and the finite
source-excluded boxes remove positive energy without
changing their lengths.  The target lifted-path
envelope similarly contains \emph{every} marked parent
leaf.  The exact angle/AGM checks alone do not verify
these global inclusions.

\item[(H4) Peripheral-cover comparison.]
The universal-cover ideal interval lifts for the
four-hole KKL canonical WAD have the stated
disjoint-translate property, and the physical
Joukowski annuli compare to the corresponding
peripheral cyclic covers with the claimed transverse
width convention.  The finite ellipses themselves
are checked algebraically, but their interpretation
as bounds on \emph{all} ideal isotopy classes
uses this analytic input.
\end{description}
\end{definition}

These hypotheses are \emph{not} inferred from passing
Python/Lean checks or from prior in-repository audit
claims.  The named classical results provide plausible
tools to discharge them, but the referee has not verified
their complete applications in the manuscript.
Accordingly:
\begin{itemize}
\item The local root and Rouch\'e statements retain
their explicit, piecewise scope and written arithmetic/analytic
argument; their implementations are independently
\emph{executable} rather than machine proofs of the tail.
\item The complete physical marking conclusions are
asserted \emph{only under (H1)--(H2)}; the full
unbuffered canonical sandwich and strict
independently optimized defect are asserted
\emph{only under (H1)--(H3)}.  The reciprocal arithmetic
still gives strict $1/100$ once those inputs hold.
\item The buffered exclusion is conditional on
(H1), (H2), and (H4); no inner class is decided.
Any depth word's identification with a named
straightening suffix needs an independent \emph{full}
kneading argument beyond a table of real signs.
\end{itemize}

The conditional statements are useful to localize the exact
remaining proof tasks, but must \emph{not} be advertised
as unconditional new theorems or as KKL amplification.
The original first-version assertions remain unmodified
under \texttt{papers/tip\_family\_preprint\_20261116/}.

\section{Selected parameter, orbit order and proper returns}\label{sec:parameter}
Throughout $x_j=x_j(c_p)$ means $f_{c_p}^j(0)$ in the physical coordinate.

Define $q=4^{-p}$ and
\begin{equation}\label{eq:basic}
 f_c(z)=z^2+c,\quad G_c=f_c^p,\quad
 P_p(c)=x_{6p}(c)+x_{3p}(c),\quad
 A_p(c)=\prod_{j=1}^{p-1}2x_j(c),\quad
 \tau_p(c)=A_p(c)x_p(c).
\end{equation}
For $3\le p\le12$ let $\sigma_p$ be the exact rational decimal in Table~\ref{tab:seeds}, and use the open chart
$(\sigma_p-10^{-75},\sigma_p+10^{-75})$.

For $p\ge13$ define the explicit near-minus-two chart
\begin{align}
 \mathcal{C}_p&=\{c=-2+\eps:14q<\eps<16q,\ 
 |x_p(c)|<14q,\ -2<\tau_p(c)<1/4\},\label{eq:chart}\\
 \mathcal{K}_p&=\{c\in\mathcal{C}_p:x_1<0,\ x_j>0\ (2\le j\le p),\
 (\operatorname{sgn}x_{jp})_{j=1}^6=(+,-,-,+,-,+)\}.\label{eq:knead}
\end{align}
These conditions contain no critical-orbit equation.
\begin{theorem}[Selected real tips]\label{thm:real}
There is exactly one simple zero $c_p$ of $P_p$ in the finite chart for $3\le p\le12$ or in $\mathcal K_p$ for $p\ge13$.  No larger-window uniqueness is claimed.  The dividing fixed point $\alpha_p$ is unique in $(x_{6p},x_{4p})$, and
\begin{align}
 x_{2p}&<x_{5p}<x_{3p}<0<x_{6p}<\alpha_p<x_{4p}<x_p,\label{eq:order1}\\
 -x_p&<-x_{4p}<x_{2p}<x_{5p}<-\alpha_p<x_{3p}<x_{6p}
 <\alpha_p<-x_{5p}<-x_{2p}<x_{4p}<x_p.\label{eq:order2}
\end{align}
The map $G_p$ is increasing on $[-x_p,0]$ and decreasing on $[0,x_p]$; $A_p<0$ and there is no earlier central critical return.  For every $p\ge13$,
\begin{equation}\label{eq:tip-rate}
 \left|c_p+2-\frac{3\pi^2}{2}4^{-p}\right|
 <500p\,16^{-p},
 \qquad c_p+2\sim\frac{3\pi^2}{2}4^{-p}.
\end{equation}
\end{theorem}

\begin{theorem}[Degree-two first and third returns]\label{thm:rouche}
Put $R_3=1/4$, $R_p=2^{-p}$ for $p\ge4$,
$U_p=D(0,R_p)$, and $V_p=\comp_0 G_p^{-1}(U_p)$.
For all $p\ge3$, $V_p\Subset U_p$ and
$G_p:V_p\to U_p$ is proper holomorphic of \emph{degree two},
with the full-complex-disk Rouch\'e gap $>1/25$.
For $p\ge13$ the explicit margin is
\begin{equation}\label{eq:p-margin}
 |A_p|R_p^2-|x_p|-T_p-R_p
 \ge 3/20-12q-1/128-2^{-p}>1/25,
\end{equation}
where $T_p$ bounds every term of degree at least four
in $G_p(z)-x_p-A_pz^2$ on the entire disk $|z|\le R_p$.
The \emph{different} third return, normalized by
$h_p(w)=A_pG_p(w/A_p)$, has a central proper
degree-two restriction into $D(0,3/10)$, with
normalized Rouch\'e gap $>1/8$.
The unrestrained polynomial has degree $2^p$; nothing
here gives a uniform modulus for subsequently shrunken
three-return annuli.
\end{theorem}

\subsection{Finite roots and interval-soundness}
\begin{proof}[Proof of Theorem~\ref{thm:real}]
First take $3\le p\le12$.  Define exact polynomial recurrences
\[
 X_0=Y_0=0,\quad X_{j+1}=X_j^2+c,\quad
 Y_{j+1}=2X_jY_j+1.
\]
Then $X_j=x_j(c)$ and $Y_j=x'_j(c)$ identically
over $\Q[c]$.  For each rational seed in
Table~\ref{tab:seeds}, the finite checker verifies
\begin{equation}\label{eq:finite-sign}
 P_p(\sigma_p-10^{-75})>0,\quad
 P_p(\sigma_p+10^{-75})<0,\quad
 \sup_{|c-\sigma_p|\le10^{-75}}
 (Y_{6p}(c)+Y_{3p}(c))<0.
\end{equation}
The endpoints are rational, and the derivative is evaluated
on the \emph{entire} chart by outward $2^{-512}$
interval recurrences, not by sampling.  The intermediate
value theorem and strict monotonicity prove one simple root
in each open chart.  Eighty-two further sign-separated
bisections locate it in an adjacent $2^{-256}$
rational root cell; every operation on the $X_j,Y_j$
arrays and every comparison is an integer inequality.
The resulting full-hull orbit and a two-sided
fixed-point sign bracket prove all inequalities in
\eqref{eq:order1}--\eqref{eq:order2}, $A_p<0$, and the
real monotonicity.  Thus the proof of each finite assertion
is a finite exact decision procedure with explicitly supplied
input seeds, recursion, precision and required output signs.
The complete-disk properness check is proved separately below.

For $p\ge13$ use $q=4^{-p}$, $c=-2+\eps$,
$14q\le\eps\le16q$.  The positive repelling
fixed point of $f_c$ is
\[
 \beta=\frac{1+\sqrt{9-4\eps}}2,\quad
 \lambda=2\beta=4-\delta,\quad
 \delta=\frac{4\eps}{3+\sqrt{9-4\eps}}\in(0,\eps).
\]
Setting $y_j=\beta-x_j$, exact substitution gives
\begin{equation}\label{eq:cheb}
 y_{j+1}=\lambda y_j-y_j^2,\quad
 y_2=\frac83\eps-\eps^2-
 \frac{4\eps^2}{3(3+\sqrt{9-4\eps})^2},\quad
 \frac83\eps-2\eps^2<y_2<\frac83\eps.
\end{equation}
For $j\ge2$ put $y_j=2(1-\cos\theta_j)$
on the positive small-angle branch, and put
$\Theta=2^{p-2}\theta_2$ and $\kappa=\eps/q$.
The alternating Taylor inequalities for cosine,
$\sin u\ge2u/\pi$ on $[0,\pi/2]$ and
$157/50<\pi<22/7$ imply
\begin{equation}\label{eq:angle-bds}
 \frac{\kappa}{6}(1-\eps)<\Theta^2
 <\frac{\kappa}{6}(1+3\eps),\qquad
 \frac{38}{25}<\Theta<\frac{41}{25},\qquad
 \theta_j<\frac{41}{50}\quad(2\le j<p).
\end{equation}
To see the main squared-angle comparison explicitly,
$1-3\eps/4<y_2/((8/3)\eps)<1$,
$\theta_2^2\le(\pi^2/4)y_2<7\eps$,
and $1\le\theta_2^2/y_2
\le(1-7\eps/12)^{-1}\le1+\eps$.
Multiplication by $4^{p-2}$ gives the first inequality.
In the above small-angle regime the perturbation of the
exact Chebyshev doubling map satisfies
\[
 (2-\eps)\theta_j\le\theta_{j+1}\le2\theta_j.
\]
Iteration, the $1$-Lipschitz property of cosine,
and $|\beta-2|<\eps/2$ yield
\begin{equation}\label{eq:escape}
 |\theta_p-\sqrt{\kappa/6}|\le(p+6)\eps,\qquad
 |x_p-2\cos\sqrt{\kappa/6}|\le3p\eps.
\end{equation}
In particular rational endpoint comparisons using
$38/25<\sqrt{14/6}<77/50$ and
$\sqrt{16/6}>81/50$ show
$x_p(-2+14q)>1/20$ and
$x_p(-2+16q)<-1/20$.

These opposite signs alone would not give
parameter uniqueness.  Differentiating
\eqref{eq:cheb} gives, with $Y_j=dy_j/d\eps$,
\[
 Y_{j+1}=(\lambda-2y_j)Y_j+\lambda' y_j,\quad
 r_j=\frac{\eps Y_j}{y_j},\quad
 r_{j+1}=\left(1-\frac{y_j}{\lambda-y_j}\right)r_j+
 \frac{\eps\lambda'}{\lambda-y_j}.
\]
The verified rational bounds
$r_2>99/100$, $0<y_j<7/10$ for $2\le j<p$,
$\sum y_j<21/20$,
$\sum y_j/(\lambda-y_j)<7/20$ and
$p\eps<1/300000$, together with
$\prod(1-a_j)\ge1-\sum a_j$,
prove $r_p>1/2$.  Since $\theta_p>3/2$,
$y_p>9/5$, this gives
\begin{equation}\label{eq:slope}
 \frac{dx_p}{dc}<-\frac{1}{20q}.
\end{equation}
Thus the window $|x_p|\le14q$ is
crossed exactly once by the parameter.

We next normalize the first return, including its
\emph{parameter derivative}.  Write $a_j=[z^2]f_c^j$,
with $a_1=1$ and $a_{j+1}=2x_ja_j$.
The exact cosine product identity
\[
 \prod_{j=2}^{p-1}\cos(2^{j-2}\theta_2)
       =\frac{\sin\Theta}{2^{p-2}\sin\theta_2}
\]
and bounds $1-\eps\le
2x_j/(4\cos(2^{j-2}\theta_2))\le1+p\eps$
give
\begin{equation}\label{eq:A-bds}
 A_p<0,\qquad 3/20<-A_pq<17/100.
\end{equation}
Differentiation of the $x_j,a_j$ recursions,
using $|x_j|<2$ before return, gives
$|x'_j|\le4^j/3$, $|a_j|\le4^{j-1}$,
$|a'_j|\le16^{j-1}/3$.
Consequently $\tau_p'=A'_px_p+A_px'_p
>7/(1000q^2)$ on $|x_p|\le14q$,
with $\tau_p$ traversing $[-2,1/4]$ only once.
In the subchart
$J=[-1795/1000,-1785/1000]$,
one has $0<x_p<12q$.

The full complex perturbation estimate below yields
\[
 h_\tau(w)=A_pG_p(w/A_p)=w^2+\tau+e_\tau(w),\quad
 |e_\tau(w)|<88q,\quad
 |\partial_we_\tau(w)|\le256q,\quad
 |\partial_\tau e_\tau(w)|<30000q^2
 \quad(|w|\le21/10).
\]
On $|w|\le19/10$, the spatial derivative bound is
$171q$.  For $g_\tau(w)=w^2+\tau$ let
$u_j=g_\tau^j(0)$ and $\Pi(\tau)=u_6+u_3$.
On $J$, rational evaluation gives
\[
 \Pi(-1795/1000)<-1/20,\quad
 \Pi(-1785/1000)>1/20,\quad
 \Pi'(\tau)>2,\quad
 |u_j|<9/5\quad(0\le j\le6).
\]
Set $E_0=D_0=0$ and inductively
\begin{equation}\label{eq:two-errors}
 E_{j+1}=4E_j+88q,\quad
 D_{j+1}=4D_j+25(2E_j+171q)+30000q^2.
\end{equation}
For $p=13$ (and therefore for all larger $p$)
the strict rational bounds are
$E_6+E_3<1/500$ and $D_6+D_3<1/5$.
This proves that the perturbed six/three-return
expression has opposite endpoint signs
and derivative $>2-1/5>0$ throughout $J$.
It has exactly one simple zero there.

To exclude a second root in the \emph{whole signed chart},
first the quadratic model on $[-2,0]$ is trapped in
a real invariant interval; the comparison recurrence
$E_{j+1}^{\rm wide}=(21/5)E_j^{\rm wide}+88q$
gives $E_6^{\rm wide}+E_3^{\rm wide}<1/400$.
The six rational closed nonroot segments
\[
\begin{split}
 [-2,-1.989],\ [-1.988,-1.982],\
 [-1.980,-1.894],\ [-1.892,-1.795],\\
 [-1.785,-1.756],\ [-1.753,-0.01]
\end{split}
\]
are each covered by \emph{1000} rational subintervals;
the outward interval evaluation on each entire cell
proves $|\Pi|>1/400$.  In the intervening neighborhoods
$[-1.989,-1.988]$, $[-1.982,-1.980]$ and
$[-1.894,-1.892]$, the required signed normalized
itinerary is contradicted respectively by the
fourth or fifth return.  On
$[-1.756,-1.753]$ put $v=h_\tau^3(0)$:
the bounds $|v|<1/40$ and
$|(h_\tau^3)'(w)|<1$ for all $w\in[0,v]$
give
\[
 h_\tau^6(0)+h_\tau^3(0)
 =v\left(2+\int_0^1(h_\tau^3)'(sv)\,ds\right),
\]
so a zero would force $v=0$, which is forbidden by the
strict $x_{3p}<0$ condition.
For $-0.01<\tau<0$, evenness gives
$e_\tau(w)=O(w^4)$; the order-four Schwarz
lemma on $|w|\le12/5$ yields
$|e_\tau(\tau)|<|\tau|/100$ and then
$h_\tau^2(0)=\tau(1+\tau)+e_\tau(\tau)<0$,
also forbidden.  Finally $\tau\ge0$
contradicts $x_p>0$.
These analytic exclusions and all finite rational
cells, \emph{not} an unisolated broad polynomial root,
prove uniqueness in $\mathcal K_p$.

At that root the normalized model's strict
six-return cut bounds with errors \eqref{eq:two-errors}
give \eqref{eq:order1} and
\eqref{eq:order2}.  The fixed point is bracketed
on the relevant interval with strictly negative
$h_\tau'$ and opposite endpoint signs.  The
negative factor $A_p$ reverses the normalized
real cut order, which is included in the
outward computations.  Finally the derivative
of $2\cos\sqrt{\kappa/6}$ has absolute value
$>1/10$ on $[14,16]$; with $|x_p|<12q$,
\eqref{eq:escape} and the mean-value theorem
give
$|\kappa-3\pi^2/2|<500pq$, exactly
\eqref{eq:tip-rate}.  The strict orbit enclosure
also gives the two monotone inverse sheets and
absence of an earlier central critical return.
\end{proof}

\subsection{Proper degree-two returns (not globally degree two)}
\begin{proof}[Proof of Theorem~\ref{thm:rouche}]
For $f_c^j(z)=x_j+a_jz^2+E_j(z)$,
$a_1=1$, $a_{j+1}=2x_ja_j$, and
$E_1=0$, the \emph{identity}
\[
 E_{j+1}(z)=2x_jE_j(z)+(a_jz^2+E_j(z))^2
\]
proves by induction, \emph{on the whole complex disk} $|z|\le R$,
\begin{equation}\label{eq:tail}
 T_1=0,\quad
 T_{j+1}=2|x_j|T_j+(|a_j|R^2+T_j)^2,\quad
 \sup_{|z|\le R}|E_j(z)|\le T_j.
\end{equation}
This is a complete tail, not a finite Taylor truncation.
For $p\ge13$, $|x_j|<2$ and $|a_j|\le4^{j-1}$.
If $t_j=4^jR^2\le1$, induction yields
$T_j\le t_j^2/128$: for $j<p$,
$t_j\le1/4$ and
$4/128+(1/4+1/512)^2<16/128$.
With $R=R_p=2^{-p}$, $t_p=1$ and
$T_p\le1/128$.  Thus
\eqref{eq:A-bds}, $|x_p|<12q$ and
$R_p\le2^{-13}$ imply \eqref{eq:p-margin}.
Every one of these strict comparisons reduces
to a rational inequality at the worst $p=13$.
For $3\le p\le12$ the original ten root packets
are propagated via \eqref{eq:tail}, respecting
the exceptional $R_3=1/4$, to prove
the identical $>1/25$ boundary margin.

For any $|y|\le R_p$, on $|z|=R_p$,
\[
 |G_p(z)-y-A_pz^2|
 \le|x_p|+T_p+R_p<|A_p|R_p^2.
\]
Rouch\'e's theorem
\cite[Chapter 4, \S5.2, corollary to Theorem 18]{AhlforsComplex}
counts exactly two solutions of $G_p(z)=y$
inside the disk, with none on its boundary.
Since $G_p(0)=x_p\in U_p$ and zero has
local multiplicity two, both solutions
near this critical value lie in the central
component.  By proper continuation this
component maps onto $U_p$ with degree two.
The strict boundary inequality and
Riemann--Hurwitz
\cite[\S17.14]{Forster1981} give the
central proper quadratic-like restriction,
$V_p\Subset U_p$, with a single critical point.

For the \emph{different} normalized third return
$H=h_\tau^3$, the model $g_\tau(w)=w^2+\tau$
has the \emph{exact} degree-eight polynomial
\[
 g_\tau^3(w)=(v^2+\tau)+(4\tau v)w^2
 +(2v+4\tau^2)w^4+4\tau w^6+w^8,
 \qquad v=\tau^2+\tau.
\]
For $\tau\in J$ one has the verified bounds
$7/5<v<10/7$, $|v^2+\tau|<13/50$,
$4|\tau|v>99/10$, $2v+4\tau^2<16$,
$4|\tau|<8$.  On the \emph{whole}
complex disk $|w|\le r=3/10$,
the first two intermediate images remain in
$|w|<21/10$, so the three-step perturbation
is bounded by $1848q$.  The Rouch\'e gap is
therefore strictly larger than
\[
 \frac{99}{10}r^2-\left(
   \frac{13}{50}+16r^4+8r^6+r^8+
   1848\cdot4^{-13}+r\right)>\frac18.
\]
The third-return central local degree is two;
both intermediate model images stay away from
additional critical points.  Rouch\'e and
Riemann--Hurwitz prove the proper central
$H$ restriction.  For $3\le p\le12$ the
separate full-disk $3p$-return check applies
\eqref{eq:tail} at $R=r/|A_p|$ with outward
$2^{-192}$ tail rounding and certifies the
same $>1/8$ strict gap and intermediate
component containment.
The three phase compact real intervals are
disjoint by \eqref{eq:order1}; shrinking
their analytic neighborhoods gives a
primitive three-return restriction
inside the original $G_p$ domain.  This
shrinking proves \emph{existence} but not
a uniform modulus of the chosen neighborhoods.
\end{proof}

\begin{table}[htbp]
\centering
\caption{Exact terminating decimal rational seeds $\sigma_p$ for $3\le p\le12$.
The selected chart has half-width $10^{-75}$; no decimal is
assumed to equal a root.  For each row, concatenate the two successive text strings without spaces to recover the exact decimal. The complete interval proof checks
endpoint signs and a negative derivative throughout that chart.}
\label{tab:seeds}
\begingroup\fontsize{7.0}{8.4}\selectfont
\begin{tabular}{@{}r l@{}}
\toprule
$p$ & Exact rational $\sigma_p$\\ \midrule
3&\shortstack[l]{\texttt{-1.786510544544324757385260673471213673058457671889172}\\\texttt{800031923968334629798587697269507048968639156733}}\\
4&\shortstack[l]{\texttt{-1.942554171494761905319728559547351970316877794253756}\\\texttt{256973061644721007908220533846556129430732097156}}\\
5&\shortstack[l]{\texttt{-1.985528169705348810408255730556821748277263841949252}\\\texttt{553593917274749073066379677891609603820336856318}}\\
6&\shortstack[l]{\texttt{-1.996382511718412027950049688816515039904858979750298}\\\texttt{286593034926678792080923860935932139817781617068}}\\
7&\shortstack[l]{\texttt{-1.999096078309175980023599670405627474503150312022034}\\\texttt{53414490734530886689705706443584075484634242529}}\\
8&\shortstack[l]{\texttt{-1.999774073397058750869233413315687996855142520380809}\\\texttt{74798691863676382684508119126543885186585434088}}\\
9&\shortstack[l]{\texttt{-1.999943523308791022549753910827383353953529788170904}\\\texttt{99897525072423167823633706849247069259463644468}}\\
10&\shortstack[l]{\texttt{-1.999985881236821742548433835039758405032921437752994}\\\texttt{01058889825653946965600383848213765034339890565}}\\
11&\shortstack[l]{\texttt{-1.999996470341035272735632678015004938386461804580811}\\\texttt{692099169946758788547010051826504894628853093167}}\\
12&\shortstack[l]{\texttt{-1.999999117587637481776828101553703232713431093876163}\\\texttt{616822020428055010613627327045591565019864984472}}\\
\bottomrule
\end{tabular}
\endgroup
\end{table}


\section{The actual little Julia arcs and KKL ideal boundaries}
\label{sec:marking}
All markings in this section refer to the \emph{physical}
slit complements of the proper central restriction.
They are not a substitute for a marked ambient puzzle
or a specification of external ray angles.

\begin{theorem}[Physical period-three marking, conditional on (H1)--(H2)]
\label{thm:marking}
Assume the geometric renormalization and physical ideal-marking
bridges \emph{(H1)}--\emph{(H2)} of
Definition~\ref{def:bridges}.
Put $t=x_{6p}=-x_{3p}>0$, and
\begin{equation}\label{eq:tree}
 K_0=[-t,t],\quad K_1=[x_{4p},x_p],\quad
 K_2=[x_{2p},x_{5p}],\quad
 \T_p=[x_{2p},x_p],\quad
 \T'_p=[-x_p,x_p].
\end{equation}
The renormalization of $G_p$ of period three is
\emph{prime and primitive}, and its \emph{actual}
little filled Julia sets are $K_0,K_1,K_2$.
The KKL Hubbard tree is exactly $\T_p$, with
dividing fixed point $\alpha_p$.
Moreover $\T'_p=G_p^{-1}(\T_p)\cap V_p$.

The target marked surface $\Om_p=U_p\setminus\T_p$
has ten consecutive \emph{inner} ideal atoms
$I_0,\ldots,I_9$, indexed by the five physical
increasing cuts
\[
 (a,b,d,e,h)=(x_{5p},x_{3p},x_{6p},\alpha_p,x_{4p}).
\]
The source $\Om'_p=V_p\setminus\T'_p$
has twenty consecutive inner ideal atoms
$S_0,\ldots,S_{19}$ indexed by
\[
 (-x_{4p},x_{2p},x_{5p},-\alpha_p,x_{3p},
 x_{6p},\alpha_p,-x_{5p},-x_{2p},x_{4p}).
\]
The two \emph{outer} ideal circles each have one
whole-component atom, not ten or twenty.
The unique new source-tree attachment is
$\T'_p\setminus\T_p=[-x_p,x_{2p})$, with
consecutive source sides $S_{19},S_0,S_1$.
\end{theorem}

\begin{proof}
The exact equation $x_{6p}=-x_{3p}$,
evenness of $G_p$ and of its iterates,
and \eqref{eq:order1} give
\begin{equation}\label{eq:portrait}
 G_p^3(0)=-t,\quad G_p^3(-t)=G_p^3(t)=t,\qquad
 G_p(K_0)=K_1,\quad G_p(K_1)=K_2,\quad
 G_p(K_2)=K_0.
\end{equation}
The full third-return disk in Theorem~\ref{thm:rouche}
has a central proper degree-two restriction.
In the normalized coordinate $w=A_pz$ its
connected filled Julia set contains the actual
real invariant interval $[-|A_p|t,|A_p|t]$:
the interval is forward invariant by \eqref{eq:portrait},
lies in the certified central disk component and
contains the sole critical point.  The Douady--Hubbard
connectedness criterion
\cite[Proposition 2]{DH1985} therefore makes the
filled Julia set connected.  By the
\emph{Straightening Theorem}
\cite[Theorem 1]{DH1985} it is hybrid-conjugate
to a unique monic centered quadratic $z^2+\gamma$.
Its critical portrait, in the normalized orientation,
is $0\mapsto u\mapsto -u\mapsto -u$ with $u>0$.
For a monic quadratic the relation
$f_\gamma^3(0)=f_\gamma^2(0)$ reduces
\emph{exactly} to $\gamma^3(\gamma+2)=0$.
The nonfixed critical point excludes $\gamma=0$,
so $\gamma=-2$.  Since $K(z^2-2)=[-2,2]$,
the filled Julia set of the third return is an arc.
It contains the entire straight segment joining
its endpoints $\pm u$, and an arc cannot contain
a proper different arc with the \emph{same}
endpoints.  Thus it is exactly this segment,
not merely homeomorphic to it.  Rescaling and
\eqref{eq:portrait} yield $K_0,K_1,K_2$ as
the three actual little filled Julia intervals.

The strict separations in \eqref{eq:order1}
make the three little filled sets pairwise
disjoint; this is \emph{primitive}
in \cite[\S2.3]{KKL2026}.  We prove primeness
rather than infer it from the word primitive.
The global polynomial $f_{c_p}$ is postcritically
finite and its critical point is strictly
preperiodic: $f^{9p}(0)=f^{6p}(0)=t\ne0$,
whereas $f^{3p}(0)=-t\ne0$.
The Misiurewicz local-connectivity/dendrite and
periodic-repulsion theorem
\cite[Theorem 27.4]{LyubichNotes}
and the subhyperbolic treatment
\cite[\S17]{Milnor2006} make its connected
Julia set a locally connected dendrite,
with all periodic cycles repelling.
The \emph{global} filled Julia sets agree,
$K(G_p)=K(f_{c_p})$, since $G_p=f_{c_p}^p$;
they are connected because the \emph{unique
critical point of the quadratic polynomial
$f_{c_p}$} has bounded orbit.  This does
not use the false premise that the
degree-$2^p$ global polynomial $G_p$
has only one critical point.  The relevant
central little filled set, once its
connectedness is proved for the restricted
proper map, is a subcontinuum of the
ambient dendrite and hence is a dendrite.

Suppose a connected period-two degree-two
renormalization of $G_p$ had little filled set
$L\ni0$.  Since the even-return critical orbit
lies in $L$, both $x_{2p}$ and $x_{4p}$
belong to it.  Uniqueness of dendrite arcs
then forces $[x_{2p},x_{4p}]\subset L$.
By the real monotonicity and \eqref{eq:portrait},
$L\cap G_p(L)$ contains the nondegenerate
interval $[x_{5p},x_{4p}]$.
Under the \emph{nestedness/straightening
transport} clause of (H1), these are
little filled Julia sets to which
McMullen's \emph{almost-disjoint little Julia
sets theorem} \cite[Theorem 7.3]{McMullen1994}
applies, so distinct such sets meet in
at most one repelling periodic point;
hence $L=G_p(L)$.  The theorem is
\emph{not} applied directly to the
unrestricted degree-$2^p$ polynomial.
But $0\in G_p(L)$ then has an inverse
$y\in L$, $y\ne0$ because $G_p(0)=x_p>0$.
The proper-domain containment demanded
by (H1) places both $0$ and $y$ in
the alleged degree-two return domain;
there they are distinct critical points
of $G_p^2$, contradicting its degree.
Without that domain containment the
original argument would be incomplete.
Thus no period-two renormalization exists.
The period-three renormalization is prime.
(The period-one case is the original $G_p$.)

Real invariance and the order \eqref{eq:order1}
show $\T_p=[x_{2p},x_p]\subset K(G_p)$ and
$G_p(\T_p)\subset\T_p$.  The three little
Julia intervals have extreme endpoints
$x_{2p}$ and $x_p$.  By dendrite arc
uniqueness any connected subset containing
them contains their real joining interval.
KKL's \emph{smallest connected Hubbard
continuum} definition
\cite[\S2.5]{KKL2026} therefore gives
precisely $\T_p$, not a conjectural drawn
subtree.  The unique fixed point
$\alpha_p\in(x_{6p},x_{4p})$ has the
dividing property by real cut order, as
in KKL \S2.4.  Finally
$G_p([-x_p,x_p])=\T_p$
under the two strictly monotone sheets.
The proper degree-two count exhausts
the two central inverse sheets; therefore
$G_p^{-1}(\T_p)\cap V_p=[-x_p,x_p]$,
not an additional off-real preimage arc.

For the \emph{ideal boundaries}, $U_p\setminus\T_p$
is a doubly connected domain with a round
outer boundary and a locally connected real
slit inner boundary.  So is
$V_p\setminus\T_p'$: the central proper degree-two preimage
component has a finite analytic
lemniscate boundary and hence locally
connected boundary (possible critical
boundary values require care);
a Jordan boundary is \emph{not}
inferred merely from the Rouch\'e count.
Also $\T'_p\Subset V_p$.
Cut either annulus along a short crosscut
joining its two boundary components.
The remaining simply connected slit
domain has locally connected boundary.
The Carath\'eodory prime-end
compactification \cite{Caratheodory1913}
provides an ideal circle for each
doubly-connected slit complement;
the further \emph{physical identification}
of all accesses with the claimed KKL
labels is precisely condition (H2).
Gluing the two crosscut banks then recovers the \emph{two}
ideal boundary circles of KKL Appendix A.1.
An interior point of the real slit has
exactly two local approach sectors
(above and below the slit) in a small
disk, and a terminal slit tip has
exactly one.  Shrinking crosscuts
show no third access is hidden.
The five inner target cut positions
therefore give ten inner ideal atoms,
and the ten source cut positions give
twenty.  The \emph{outer} circle
is entirely uncut, hence contributes
one ideal atom.  The physical order
\eqref{eq:order2} identifies all source
cuts.  These are actual
Carath\'eodory/KKL slit-side atoms,
without an unproved external ray-angle
atlas or an extra family-dependent
marking hypothesis.  Since $\T'_p$
extends $\T_p$ only to the left,
the set difference is the asserted
\emph{single} attachment, with
sides $S_{19},S_0,S_1$ by cyclic order.
\end{proof}

\begin{theorem}[Covering, immersion and incidence, conditional on (H1)--(H2)]
\label{thm:incidence}
Assume \emph{(H1)}--\emph{(H2)} of
Definition~\ref{def:bridges}, including the KKL
prime-end and pulled-marking compatibility.
On $\Om'_p$ the proper unbranched map
$G_p:\Om'_p\to\Om_p$ induces the
\emph{cover} $G_{p*}$, while the
nonproper inclusion $\iota:\Om'_p\hookrightarrow
\Om_p$ induces the \emph{set-valued proper-path
restriction} $\iota^*$.  With the above cyclic
indices their exact distinct tables are
\begin{equation}\label{eq:cover}
 G_{p*}(S_s)=I_{s\bmod10},\quad 0\le s<20,
\end{equation}
\begin{equation}\label{eq:blocks}
\begin{array}{c|l@{\qquad}c|l}
I_j&\iota^*I_j&I_j&\iota^*I_j\\ \hline
I_0&S_{18},S_{19},S_0,S_1,S_2&
I_5&S_{10}\\
I_1&S_3,S_4&I_6&S_{11},S_{12},S_{13}\\
I_2&S_5&I_7&S_{14}\\
I_3&S_6&I_8&S_{15}\\
I_4&S_7,S_8,S_9&I_9&S_{16},S_{17}
\end{array}.
\end{equation}
The attachment makes the $I_0$ block
not fully proper.  For the specified
crossing-attachment relative isotopy
class the two proper source pairs are
$(S_3,S_1)$ and $(S_{19},S_{17})$;
for the avoiding class the one pair
is $(S_3,S_{17})$.  Their inclusive
short source spans are $3,3,7$;
their \emph{cover-pushed} target
pairs are $(I_3,I_1),(I_9,I_7),(I_3,I_7)$,
with short spans $3,3,5$.
All pass the KKL A.1 filter $>2$.
No arbitrary A.1 parent is asserted
to have two restrictions.

For $E_0=(x_{5p},x_{3p})$,
$E_1=(x_{6p},\alpha_p)$,
$E_2=(\alpha_p,x_{4p})$,
the image incidence, indexed by output
rows and input columns, is
\begin{equation}\label{eq:M}
 M=\begin{pmatrix}0&0&1\\1&0&1\\1&1&0\end{pmatrix},
 \quad\det(zI-M)=(z+1)(z^2-z-1),\quad
 \rho(M)=\varphi=\frac{1+\sqrt5}{2}.
\end{equation}
The positive left Perron vector is
$(\varphi,1,\varphi)$.
Target zero-old-edge atom types are
$I_0,I_2,I_5,I_8$.
The atomic parent crest $I_0$ and
the cover-pushed first-child crests
$I_2,I_8$ have no old-edge carrier.
The two banks $I_1,I_9$ are \emph{one}
physical edge $E_0$, not two edges.
\end{theorem}
\begin{proof}
The restriction $G_p$ has no
critical points in $\Om'_p$:
its only central critical point
$0$ was removed with $\T'_p$.
The degree-two properness of
Theorem~\ref{thm:rouche} makes
$G_p:\Om'_p\to\Om_p$ an unbranched
cover.  KKL Appendix A.3.1
extends it to the ideal circle,
where one rotation-compatible
sheet labeling gives
\eqref{eq:cover}; every target
atom has two preimages.

By contrast the \emph{inclusion}
$\Om'_p\hookrightarrow\Om_p$
is nonproper precisely because
the source slit extends to the
left of the target one.  KKL
\cite[Proposition 6.5]{KKL2026}
describes the pullback of an
ideal interval as inverse-side
accesses \emph{together with}
Carath\'eodory boundaries of
attached tree components.
Traversing the physical
ten source cuts in
\eqref{eq:order2} and inserting
the three attachment sides
produces successively the
ten consecutive source blocks
\eqref{eq:blocks}, which
partition all twenty indices
exactly once.  The attached
part of the $I_0$ block
maps by inclusion to target
\emph{interior}; no point map
$I_0\mapsto S_s$ can replace
$\iota^*$ here.  KKL
\cite[Proposition 6.6 and
Appendix A.3.2--A.3.4]{KKL2026}
then gives the proper-path
restriction semantics.

A crossing arc whose right
intersection belongs to
$(x_{5p},-\alpha_p)$ and
whose left intersection belongs
to $(-x_{4p},x_{2p})$ splits
at the attached tree into two
proper source components, whose
upper/lower endpoints, read from
the cut order, are $(3,1)$ and
$(19,17)$.  If the specified
comparison arc avoids the attached
piece and its left endpoint is
left of $-x_p$, it leaves the
one proper pair $(3,17)$.
For twenty cyclic atoms the
inclusive short lengths are,
respectively, $3$, $3$ and $7$;
for example the shorter arc
from $S_3$ to $S_{17}$ is
$S_{17},S_{18},S_{19},S_0,S_1,S_2,S_3$.
Applying the \emph{different}
cover table gives target spans
$3,3,5$.  Every one exceeds two,
so none is removed by KKL A.1.
This is a statement about these
\emph{relative isotopy classes},
not a selection theorem for all
A.1 Waves.

The five cut-end images under
$G_p$ have index sequence
$(2,4,4,3,0)$, by the exact
orbit identity and fixed
point.  On the old real edges
$(a,b),(d,e),(e,h)$,
coverage counts give:
$E_0$ maps across $E_1,E_2$,
$E_1$ across $E_2$,
and $E_2$ across $E_0,E_1$.
This is exactly \eqref{eq:M}.
Its characteristic factorization
is a direct determinant.
Since $M^4$ has strictly
positive entries, Perron--Frobenius
\cite[Chapter XIII, \S2]{Gantmacher1959}
identifies the positive eigenvalue
$\varphi$ as the simple spectral
radius; alternatively
$(\varphi,1,\varphi)M
=\varphi(\varphi,1,\varphi)$
is a direct check.  The
real-edge carriers among the
ten cyclic banks are found
by looking at which atom
interiors lie over the three
old open intervals; the
remaining isolated atomic
types are $0,2,5,8$.
The given closed atomic crests
avoid these carriers.  This
does not assert an edge-free
crest for every fine
subsector or later iterate.
\end{proof}

\section{Conditional canonical marked widths and unbuffered defect}
\label{sec:wave-proofs}
The word \emph{canonical} here describes the
specified base/crest \emph{unbuffered A.1 Wave class}.
It does not identify KKL's distinct
buffered canonical WAD of Appendix A.2.3.

\begin{definition}[Extremal width and three independent suprema]
For a family $\Gamma$ of locally rectifiable proper
curves in a marked Riemann surface $Y$, set
\begin{equation}\label{eq:width}
 \W_Y(\Gamma)=
 \inf\left\{\int_Y\rho^2\,dA:
    \rho\ge0,\ \int_\gamma\rho\,|dz|\ge1
       \text{ for every }\gamma\in\Gamma\right\}.
\end{equation}
If $Y=(0,L)\times(0,H)$ and $\Gamma$ consists
of the vertical leaves, $\W_Y(\Gamma)=L/H$.
Put $P_p=\sup\W(\mathcal S)$ over \emph{all}
KKL A.1 Waves on $\Om_p=U_p\setminus\T_p$
with base $(I_9,I_0,I_1)$ and crest
containing $I_0$.  Independently put
$C_p^+$ and $C_p^-$ equal to the suprema
of all A.1 Waves on
$\Om'_p=V_p\setminus\T'_p$
with bases $(S_1,S_2,S_3)$ and
$(S_{17},S_{18},S_{19})$, and crests
$S_2$ and $S_{18}$, respectively.
Reflection across $\R$ identifies
the two child classes isometrically,
hence $C_p^+=C_p^-=C_p$.
No attainment, selection of a
common parent, or bijection of
all arbitrary parent and child
families is presumed.
\end{definition}

\begin{lemma}[Marked-closure identity under (H2)]\label{lem:M}
Assume the complete KKL physical ideal-marking
compatibility \emph{(H2)} in
Definition~\ref{def:bridges}.
The unbuffered marked surfaces used in the
canonical-sandwich and strict-gain calculations
are \emph{identically} $\Om_p$ and $\Om'_p$ as
embedded complex plane domains.  By
Theorems~\ref{thm:marking} and \ref{thm:incidence},
the identity extends to their physically
marked KKL Carath\'eodory ideal closures.
It therefore bijects \emph{every} A.1
proper-path class and Wave, with base,
ends and crest, and preserves its conformal
extremal width.  Thus the earlier
family-specific identification hypothesis
$M$ is discharged \emph{on these two
selected physical slits}.  No analogous
claim is made for a different physical
puzzle surface or for the four-hole WAD.
\end{lemma}
\begin{proof}
The point sets and complex coordinates
are exactly $U_p\setminus[x_{2p},x_p]$
and $V_p\setminus[-x_p,x_p]$, not
homeomorphic substitutes.  The
Carath\'eodory identification in
Theorem~\ref{thm:marking}
matches every marked side and every
outer whole-component circle.
KKL Appendix A.1 defines properness,
the forbidden length-at-most-two
subordination, relative-endpoint
homotopy into \emph{every} ideal collar,
simple laminations, their smallest
base multi-interval and its crest.
All these conditions are invariant
under the \emph{same marked-closure
identity} in both directions.
For a conformal coordinate map
$w=\phi(z)$ an admissible density
transforms as
$\rho_w(\phi(z))=\rho_z(z)/|\phi'(z)|$;
the curve lengths and area integrals
in \eqref{eq:width} are unchanged.
The identification here is the identity
and the auxiliary positive rescaling
$\xi=-A_pz$ is conformal.
Thus both classes and their
\emph{actual} widths are equal.
A bare topological equivalence
would not give the width equality.
\end{proof}

\begin{theorem}[Conditional canonical marked Wave sandwich]
\label{thm:sandwich}
Assume \emph{(H1)}--\emph{(H3)} of
Definition~\ref{def:bridges}.
For every $p\ge3$ put
\[
 m_p=\frac{x_{2p}+x_{5p}}2,\quad
 r_{0,p}=\frac{x_{5p}-x_{2p}}2,\quad
 r_{1,p}=\min\{-\alpha_p-m_p,\ x_{4p}+m_p\},\quad
 W_p^{\rm rad}=\frac1{2\pi}\log\frac{r_{1,p}}{r_{0,p}}.
\]
Write $P_{p,2}=P_{p,\pp}$ for the smaller parent
supremum requiring that \emph{every}
leaf has both named surviving proper
A.3.2 children \emph{in their designated
source classes}, with the inclusion
injective on the union of the two
child laminations.  For the full lifted-path
envelopes $Q_p,Q(L_{+,p}),Q_{\rm plane,p}$
defined in the proof,
\begin{equation}\label{eq:sandwich}
 W_p^{\rm rad}\le P_{p,2}\le P_p
 \le Q_p\le Q(L_{+,p})\le Q_{\rm plane,p}.
\end{equation}
For $p\ge13$,
\[
 0.26<P_p<0.78,\qquad
 0.26<\liminf_{p\to\infty}P_p
 \le\limsup_{p\to\infty}P_p<0.78.
\]
The limiting inequalities \emph{do not}
assert convergence or equality to an
extremal quadrilateral.
\end{theorem}
\begin{proof}
By \eqref{eq:order1}--\eqref{eq:order2},
$0<r_0<r_1$ (we suppress $p$).
For $r_0<s<r_1$ consider the open
circle $\gamma_s(t)=m+s e^{2\pi it}$,
$0<t<1$.  Its real right meeting
$q_s=m+s$ satisfies
$x_{5p}<q_s<-\alpha_p<x_{3p}$,
while its left meeting $w_s=m-s$
lies in $(-x_{4p},x_{2p})$.
Thus the \emph{only} removed target
slit point is the right meeting,
whose two accesses are precisely
$I_9,I_1$.  The curve goes around
the \emph{left} tip, so its
relative-endpoint class is
subordinate to the three-atom
arc $(I_9,I_0,I_1)$ and cannot
be subordinate to any one-
or two-atom arc.  This checks
KKL A.1 properness and the crest,
rather than assuming every
circle gives a Wave.

The complete closed radial gate disk
lies in the \emph{central} $V_p$.
For $3\le p\le12$ the whole-disk
bound and real connector are
certified by the finite Rouch\'e
interval checker; it is not a
test of one radius.  For $p\ge13$,
the normalized $h_p(w)=w^2+\tau_p+e_p(w)$
bounds the entire gate by
$|A_pz|<2$, and
\[
 |G_p(z)|<\frac7{|A_p|}
 <\frac{140}{3}4^{-p}
 <\frac35\,2^{-p}<R_p.
\]
Since the gate center belongs to
the certified central real pullback,
connectedness of the gate disk
places the whole disk in $V_p$.
The same circle in $\Om'_p$
is cut by the \emph{one} newly
attached real interval into two
proper open semicircles.  Their
source endpoints are $(S_3,S_1)$
and $(S_{19},S_{17})$, so both
survive the strict $>2$ atom
filter and meet distinct
immersion blocks.

The log chart
$z\mapsto\log((z-m_p)/r_{0,p})$
identifies the cut collar with
$(0,\log(r_1/r_0))\times(0,2\pi)$.
The vertical-leaf extremal density
and Cauchy--Schwarz give width
$\log(r_1/r_0)/(2\pi)$.
The two semicircle children have
angular height $\pi$, so each has
\emph{exactly} twice that width.
This is a width identity for
\emph{these chosen leaves only}.
The A.1 family inclusion and
two A.3.2 survivals prove
$W_p^{\rm rad}\le P_{p,2}\le P_p$.

To prove the upper bound for
\emph{every} parent leaf rather
than only circular ones, use the
Joukowski uniformization of the
plane complement of $\T_p$.
Set $m=(x_{2p}+x_p)/2$,
$r=(x_p-x_{2p})/2$ and
\[
 J(\zeta)=m+\frac r2(\zeta+\zeta^{-1}),
 \quad \zeta=-e^{s+it},\qquad
 a=\arccos\frac{m-x_{5p}}r,\quad
 b=\arccos\frac{m-x_{3p}}r,\quad0<a<b<\pi.
\]
The left-tip upper/lower
end-interval lifts in the
logarithmic cover are
$E_-=\{s=0,-b<t<-a\}$
and $E_+=\{s=0,a<t<b\}$.
Let $\widetilde\Gamma_p$
be \emph{all} proper
locally rectifiable curves
in the lifted physical
target joining these
sets.  KKL A.1
subordination to the
specified \emph{left-tip}
base selects precisely this
lift: a route around the
right tip, an additional
whole-slit winding, or
an outer-boundary endpoint
does not belong to that
relative-endpoint class.
Every admissible parent
leaf therefore lifts into
$\widetilde\Gamma_p$.
No circularity, simplicity
or child property is assumed
of $\widetilde\Gamma_p$.

One cannot equate conformal
width upstairs and downstairs
under a noninjective cover.
Here is the required
inequality proof.
For an admissible finite-area
density $\widetilde\rho$
on the lifted surface, define
on the physical surface
\begin{equation}\label{eq:period}
 \rho(\pi z)^2=
 \sum_{n\in\mathbb Z}
 \widetilde\rho(z+2\pi i n)^2.
\end{equation}
The sum is well-defined
almost everywhere; Tonelli
gives equal total areas,
and pointwise $\rho$ is at
least the chosen lifted
density along any selected
lift of a parent leaf.
Thus the pushdown density
is admissible for every
parent leaf.  Taking
infima proves
$P_p\le\W_{\widetilde\Om_p}
(\widetilde\Gamma_p)=:Q_p$.
This is an inequality, not
a width identity under
a covering.

The lifted physical
domain lies in the
upper full strip envelope
$0<s<L_+$, where
\[
 B_+=\frac{R_p+m}{r},\qquad
 \rho_+=B_++\sqrt{B_+^2+1},\qquad
 L_+=\log\rho_+.
\]
Indeed $|J(\zeta)|\ge
(r/2)(|\zeta|-|\zeta|^{-1})-m
>R_p$ when $|\zeta|>\rho_+$.
Conformal monotonicity
for the \emph{full lifted
path class}, followed by
exponentiating the strip
and applying the
Schwarz--Christoffel
quadrilateral formula
\cite[Chapter I, \S\S1--4]{AhlforsInvariants},
gives
\begin{align}
 Q_p&\le Q(L_+),&
 Q(L)&=\frac{K(k_L)}{K(\sqrt{1-k_L^2})},&
 k_L&=\frac{\sinh(\pi(b-a)/(2L))}
             {\sinh(\pi(b+a)/(2L))},\label{eq:Q}\\
 Q(L_+)&\le Q_{\rm plane,p},&
 Q_{\rm plane,p}&=\frac{K(k_\infty)}
                       {K(\sqrt{1-k_\infty^2})},&
 k_\infty&=\frac{b-a}{b+a}.\label{eq:Qplane}
\end{align}
Here $K(k)=\int_0^{\pi/2}
(1-k^2\sin^2\theta)^{-1/2}\,d\theta$;
the last inequality uses
$\sinh u/\sinh v\le u/v$
for $0<u<v$ and
monotonicity of $K(k)/K(k')$.
Equations \eqref{eq:Q}--\eqref{eq:Qplane}
are \emph{full all-path
upper envelopes}, not
claimed optimizer widths.

For $3\le p\le12$
the strict rational certificates
give $W_p^{\rm rad}>0.25$
and $Q(L_+)<0.78$,
with the individual finite
rows in the source audit.
For $13\le p\le16$
independent adjacent-dyadic
root cells give
$W_p^{\rm rad}>0.2604925$
and $Q(L_+)<0.771$.
For all $p\ge17$
the normalized-coordinate
error estimates yield
\begin{equation}\label{eq:tail-sandwich}
 W_p^{\rm rad}>
 0.2604925-7\,000\,000\,4^{-p}>0.26,\qquad
 Q_{\rm plane,p}<0.7750748<0.78.
\end{equation}
The first inequality is
indeed $>0.26$ at $p=17$
and improves thereafter.
For the upper error box,
$|\tau_p-\tau_*|\le15246q$
and the $j=2,3,5$ cut
errors are bounded by
$(30\cdot15246+88(4^j-1)/3)q$;
the box at $p=17$
contains every later one.
The exact rational-interval
cosine, logarithm and
elliptic AGM bounds
certify the numerical
claims.  As $p\to\infty$,
\eqref{eq:tail-sandwich}
gives a lower liminf at
least $0.2604925$ and
an upper limsup at most
$0.7750748$; both are
strictly inside $(0.26,0.78)$.
No continuity theorem
for the unrestricted
supremum has been used.
\end{proof}

\begin{theorem}[Strict \emph{conditional} gain of independent
unbuffered Wave suprema]\label{thm:gain}
Assume \emph{(H1)}--\emph{(H3)} of
Definition~\ref{def:bridges}.
For every integer $p\ge3$,
$P_p,C_p^+,C_p^-$ are
positive and finite and
\begin{equation}\label{eq:gain}
 D_p:=\frac1{C_p^+}+\frac1{C_p^-}
             -\frac1{P_p}
   =\frac2{C_p}-\frac1{P_p}>\frac1{100}.
\end{equation}
The strict arithmetic cases
are
\begin{align}
 P_3&>49/100,& C_3&<487/500,&
 D_3&>300/23863>1/100,\nonumber\\
 P_4&>511/1000,& C_4&<197/200,&
 D_4&>7400/100667>7/100,\label{eq:gain-cases}\\
 P_p&>581/1000,& C_p&<247/250,&
 D_p&>43500/143507>3/10\quad (5\le p\le12),
 \nonumber\\
 P_p&>13/20,& C_p&<63/50,&
 D_p&>40/819>1/100\quad (p\ge13).\nonumber
\end{align}
These are strict inequalities for the
\emph{full} marked A.1 suprema;
they do not assume the suprema
are attained or paired.
\end{theorem}
\begin{proof}
Use the positive conformal
normalization $\xi=-A_pz$,
whose scale is $\lambda_p=-A_p>0$
(and distinguish it from
the negative-order normalization
used in the real recurrence).
Write $X_j=\lambda_px_{jp}$,
$\rho_p=\lambda_pR_p$,
$m=(X_2+X_1)/2$ and
$r=(X_1-X_2)/2$.  In the
target Joukowski logarithmic
coordinates the conformal map
to the physical target is
\[
 \mathcal{J}_P(-e^{s+it})=
 m-r\cosh(s+it).
\]
It is injective for $s>0$,
$|t|<\pi$; it meets the
left-tip electrodes at the
same two angles $0<a<b<\pi$
as in the sandwich proof.
We construct an \emph{embedded
full quadrilateral foliation}
of A.1 Waves for a lower
bound, not merely evaluate an
abstract all-path modulus.

For $p=3,4$ take the right
circular segment
$Q_{h,T}=\{\Re z>0,\ |z+d|<L\}$,
with $T=31/10$, with
$h=83/100$ or $88/100$
respectively, and
\[
 d=\frac{T^2-h^2}{2h},\qquad
 L=\frac{T^2+h^2}{2h}.
\]
The ratio $(z-iT)/(z+iT)$
maps it to a wedge of opening
$2\arctan(h/T)$; raising to
$\eta=\pi/(2\arctan(h/T))$
maps to a half-plane.
On the imaginary diameter
the two electrode values
are expressed by
$u=\eta\operatorname{artanh}(a/T)$,
$v=\eta\operatorname{artanh}(b/T)$.
The symmetric cross-ratio
modulus is
\begin{equation}\label{eq:segmod}
 k=\frac{\sinh(v-u)}{\sinh(v+u)},\qquad
 W_{\rm seg}=\frac{K(k)}{K(\sqrt{1-k^2})}.
\end{equation}
The full-domain inequalities
$b<T<\pi$ and
$|m|+r\cosh h<\rho_p$
are certified over the
selected \emph{entire}
complex quadrilateral,
so its SC rectangle
\emph{foliation} maps
to a family of simple
disjoint left-tip A.1
paths with base
$(I_9,I_0,I_1)$.
The exact interval
AGM evaluation
of \eqref{eq:segmod}
proves $P_3>49/100$
and $P_4>511/1000$.
A right semidisk of radius
$5/2$ would \emph{not}
satisfy the $p=3,4$
physical containment
and is not used.

For $5\le p\le12$
take the right semidisk
$\{\Re z>0,\ |z|<R\}$
with $R=7/5$.
The Möbius map
$z\mapsto2Rz/(R^2-z^2)$
takes it to a half-plane;
on the imaginary diameter
the real symmetric
electrodes have values
$B_R(t)=2Rt/(R^2+t^2)$.
The quadrilateral ratio
and embedded Wave width are
\begin{equation}\label{eq:semidiskmod}
 k=\frac{B_R(b)-B_R(a)}{B_R(b)+B_R(a)},
 \qquad W_{\rm semi}=\frac{K(k)}
                              {K(\sqrt{1-k^2})}.
\end{equation}
The \emph{whole} semidisk
is inside the target
because $b<R<\pi$
and $|m|+r\cosh R<\rho_p$
on each exact finite chart.
The SC rectangle foliation
therefore supplies a lower
\emph{actual A.1 parent
competitor} of width
$>581/1000$ for every
$5\le p\le12$.

For all $p\ge13$ use the
same construction with
$R=5/2<\pi$.
The largest normalized
orbit interval box is
at $p=13$; it proves
$b<R$, $|m|<2$, $r<2$,
$e^{5/2}<13$, and
\[
 |\mathcal{J}_P(-e^{s+it})|
 \le|m|+r\cosh R<16
 <\frac3{20}\,2^p
 <\lambda_pR_p.
\]
The last strict gate uses
\eqref{eq:A-bds}.
All later normalized boxes
are nested inside the $p=13$
box.  Exact evaluation of
\eqref{eq:semidiskmod}
therefore proves the
\emph{all-period} strict
embedded width $P_p>13/20$.
\emph{Correction of 2026-11-23:} the original arithmetic implementation
divided the entire sum $30\cdot15246+88(4^j-1)$ by $3$, whereas only
the second summand carries that denominator.  The corrected full
$p=13$ cut box certifies the stated $13/20$ floor.  The strict
$1/100$ defect conclusion survives with the corrected $40/819$
tail constant.
The embedding, injectivity,
correct homotopy and
physical gate are checked
independently of the
abstract SC formula.

For the \emph{upper} child
bound, let
\[
 a_c=\arccos(x_{4p}/x_p),\
 b_c=\arccos(-x_{2p}/x_p),\
 c_c=\arccos(-x_{5p}/x_p),\
 d_c=\arccos(\alpha_p/x_p).
\]
The strict cuts give
$0<a_c<b_c<c_c<d_c<\pi$.
Lift the entire source
complement of its slit
into the universal
right-half-plane chart
$\xi=x_p\cosh(s+it)$.
The \emph{full} lifted
four-electrode quadrilateral
has cross-ratio
\begin{equation}\label{eq:child}
 k_c^2=
 \frac{(b_c-a_c)(d_c-c_c)}
 {(c_c-a_c)(d_c-b_c)},\qquad
 E_p=\frac{K(k_c)}{K(\sqrt{1-k_c^2})}.
\end{equation}
By the square-sum metric
periodization \eqref{eq:period},
any admissible physical
child Wave width is at most
this \emph{all-path}
envelope.  In particular
no child maximal Wave is
identified with an
elliptic optimizer.

For $3\le p\le12$
there is a \emph{strict}
source-domain saving.
The actual source lies in
$|z|<R_p$.  In its log
Joukowski chart the open
unit box
\[
 B_p=(s_p,s_p+1)\times(0,1),\quad
 (s_p)_{p=3}^{12}=
 (2,2,2,3,4,4,5,6,6,7)
\]
lies \emph{outside} this
source because the
rational Taylor lower
bound certifies
$x_p\sinh(s_p)>R_p$,
while
$|\cosh(s+it)|^2=
 \sinh^2s+\cos^2t$.
An admissible extremal
metric for the full
half-plane rectangle
can be periodized and
then restricted to the
actual source.  It remains
admissible for every
physical child path,
but its area is smaller
by the positive energy
of this entire omitted
box.  The four-electrode
SC derivative, bounded
below on the box, and
the rectangle height
bound certify
\[
 \int_{B_p}\rho_{\rm SC}^2\,dA>
 \frac1{5\,000\,000},\qquad
 C_p\le E_p-\frac1{5\,000\,000}.
\]
This is an \emph{actual
physical source metric}
saving, not the fictitious
claim that a target-parent
semidisk fits inside
the source.  The
$2^{-512}$ outward
evaluations of
\eqref{eq:child} and
the missing energy give,
respectively,
$C_3<487/500$,
$C_4<197/200$ and
$C_p<247/250$ for
$5\le p\le12$.

For $p\ge13$ a simpler,
weaker estimate suffices:
the uniform normalized
orbit/cut enclosure
around the model root
$\tau_*$, with
$|\tau_p-\tau_*|\le15246q$,
contains every actual
child cross-ratio in
the $p=13$ worst-case
interval box.
The complete full
right-half-plane
elliptic envelope
\eqref{eq:child} is
strictly $<63/50$
there, using rational
cosine inverses and
AGM remainder bounds.
Consequently $C_p<63/50$
\emph{for every} $p\ge13$,
without any missing-box
assumption at large $p$.

Finally take reciprocals
in the \emph{correct}
direction.  From $P_p>L_p$
and $0<C_p<U_p$,
\[
 \frac2{C_p}-\frac1{P_p}>
 \frac2{U_p}-\frac1{L_p}.
\]
Substitution yields
\[
 \frac2{487/500}-\frac1{49/100}
  =\frac{300}{23863},\quad
 \frac2{197/200}-\frac1{511/1000}
  =\frac{7400}{100667},
\]
\[
 \frac2{247/250}-\frac1{581/1000}
  =\frac{43500}{143507},\quad
 \frac2{63/50}-\frac1{13/20}=\frac{40}{819}.
\]
Each is strictly greater than the
claimed threshold.  All three
suprema are defined on their
independent full classes;
no common parent Wave has
been selected in this argument.
\end{proof}

\begin{proposition}[Conditional leafwise-pairable loss]
\label{prop:loss}
Assume \emph{(H1)}--\emph{(H3)} of
Definition~\ref{def:bridges}.
Let $P_{p,\pp}$ be the
supremum over \emph{only}
those parent A.1 Waves
whose every leaf has both
named A.3.2-surviving
proper source restrictions,
for which $\iota$ is injective
on their union, and for
which the two restrictions
belong to the two
designated source classes.
Then
\[
 P_{p,\pp}\le C_p/2,\qquad
 \frac2{C_p}-\frac1{P_{p,\pp}}\le0.
\]
\end{proposition}
\begin{proof}
For each such parent
$\mathcal S$, KKL
\cite[Lemma A.11]{KKL2026}
gives the width inequality below.
Proposition A.12 is a \emph{separate}
combinatorial restriction theorem,
not the source of this inequality.
With the lemma's every-leaf and
injectivity hypotheses verified
(or imposed as here), we have, with its
\emph{properness and
injectivity hypotheses},
\[
 \W(\mathcal S)\le
 \W(\mathcal S_+)\oplus\W(\mathcal S_-),
 \qquad a\oplus b=\frac{ab}{a+b}.
\]
By definition each
source width is at most
the independently optimized
$C_p$, so their harmonic
combination is $\le C_p/2$.
Taking the supremum over
the \emph{smaller}
pairable class gives
the first inequality,
and reciprocal inversion
gives the second.
In particular
Theorem~\ref{thm:gain}
proves $P_p>P_{p,\pp}$.
It \emph{does not}
provide gain for
the two restrictions
of any one parent.
\end{proof}

\section{The buffered canonical WAD is a different object}
\label{sec:buffer}
Let $K_0,K_1,K_2$ be the three
\emph{actual} little Julia intervals in
\eqref{eq:tree}.  Under (H1)--(H2), the minimal
four-boundary KKL surface is
\begin{equation}\label{eq:four}
 Y_{K,p}=D(0,R_p)\setminus(K_0\cup K_1\cup K_2).
\end{equation}
Its four \emph{whole ideal boundary
components}, rather than ten or
twenty individual cut atoms, are
$I_{\rm out},I_{K_0},I_{K_1},I_{K_2}$.
For an A.1 admissible unoriented
simple ideal class $\gamma$,
choose its two \emph{universal-cover}
ideal interval lifts, and put
$q_\gamma$ equal to their full
quadrilateral extremal width.
KKL \cite[Appendix A.2.3--A.2.5]{KKL2026}
defines the buffered weight
$w_\gamma=(q_\gamma-2)_+$ and
$\lVert W_{\rm can}\rVert=\sum_\gamma
w_\gamma$.  The exceptional
unbuffered annulus convention
does not apply to \eqref{eq:four}.
Self-endpoint types can use
distinct lifts and must not be
summarily discarded.

\begin{lemma}[Uniform cyclic-cover exclusion test]
\label{lem:peripheral}
Let $I$ be any minimally
marked nonsimply connected
peripheral ideal circle
of a finite-type hyperbolic
surface which is not the
exceptional marked annulus.
Write $\W(I)$ for the
width of its peripheral
cyclic cover.  For \emph{every}
A.1 admissible class $\gamma$
incident on $I$, including
essential self-endpoint
classes using different lifts,
\begin{equation}\label{eq:peripheral}
 q_\gamma\le
 \frac{K(\sqrt{1-e^{-\pi\W(I)}})}
 {K(e^{-\pi\W(I)/2})}.
\end{equation}
If $\pi\W(I)<\log(17+12\sqrt2)$,
then $q_\gamma<2$ and
$w_\gamma=0$ for all such
classes, with \emph{no}
winding-word cutoff.
\end{lemma}
\begin{proof}
Uniformize the universal
cover as $\mathbb H$ with a
lift of $I$ equal to
$(0,\infty)$ and its
primitive stabilizer
$\tau(z)=\lambda z$,
$\lambda>1$.  Writing
$\ell=\log\lambda$,
the logarithmic cyclic
quotient is a strip
of height $\pi$ and
circumference $\ell$.
Its transverse width is
$\ell/\pi$, not the reciprocal:
$\ell=\pi\W(I)$.

Another ideal lift
is, after normalization,
$J=(-b,-a)$ with $0<a<b$.
The translates $J$ and
$\tau J$ have disjoint
interiors: different
lifts of ideal boundary
components do not overlap.
If a translate stabilized
$J$ itself, its peripheral
stabilizer would coincide
with that of $I$,
forcing the same
boundary lift and an
A.1-boundary-subordinate
class, which is excluded.
Thus $b/a\le\lambda$.
For ideal intervals
$(0,\infty)$ and $(-b,-a)$
in $\mathbb H$, the
Schwarz--Christoffel
four-point formula gives
$k^2=1-a/b$ and
$q_\gamma=K(k)/K(\sqrt{1-k^2})$.
Since this ratio increases
strictly with $k$, we
obtain \eqref{eq:peripheral}.
The descending quadratic
Landen identity, applied
at $k_0=1/\sqrt2$, gives
\[
 \frac{K(\sqrt{12\sqrt2-16})}
      {K(\sqrt{17-12\sqrt2})}=2,
 \qquad
 1-(12\sqrt2-16)=17-12\sqrt2
  =(17+12\sqrt2)^{-1}.
\]
Hence the strict peripheral
length cutoff stated in
the lemma implies
$q_\gamma<2$.
The continuum inputs are
the uniformization,
peripheral stabilizers,
SC quadrilateral formula
and Landen's identity,
not Python theorems.
\end{proof}

\begin{theorem}[Conditional exclusion of outer-incident buffered weights
for $p\ge11$]\label{thm:buffer}
Assume \emph{(H1)}, \emph{(H2)} and
\emph{(H4)} of Definition~\ref{def:bridges}.
For $p=11,12$ and
\emph{every} integer $p\ge13$,
every KKL buffered class on
$Y_{K,p}$ incident on
$I_{\rm out}$ has $q_\gamma<2$
and $w_\gamma=0$,
including all possible
essential outer self-classes.
Moreover
\[
 0\le\lVert W_{\rm can}\rVert<619/200.
\]
Only six \emph{inner
endpoint types} remain
possible: $(K_i,K_j)$
for $0\le i\le j\le2$.
This is \emph{not}
an enumeration or exclusion
of all infinitely many
inner isotopy classes.
\end{theorem}
\begin{proof}
In the positive scale
$\xi=-A_pz$ a real segment
$[l,r]$, with midpoint
$m$ and half-length $a$,
has Joukowski chart
$\xi=m+(a/2)(w+w^{-1})$.
An embedded ellipse
$|w|=t>1$ has horizontal
semiaxis $(a/2)(t+t^{-1})$.
If its entire annulus
is contained in the
physical complement
of all four holes,
its essential cyclic
annulus lifts into the
corresponding peripheral
cover.  The explicit
potential $\log|w|/\log t$,
extended by constants,
has Dirichlet energy
$2\pi/\log t$.
The elementary annular
capacity comparison
therefore yields
\begin{equation}\label{eq:collar}
 \W(I)\le\frac{2\pi}{\log t},
 \qquad \pi\W(I)\le
 \frac{2\pi^2}{\log t}.
\end{equation}
For the \emph{outer}
component, use the whole
tree carrier
$[x_{2p},x_p]$ to build
the surrounding ellipse;
the carrier is not
misidentified as a
fourth Julia boundary.
The outward dyadic
finite chart checks give
an embedded outer ellipse
with $t=300$ at $p=11$
and $t=512$ at $p=12$,
both with
$\pi\W(I_{\rm out})<7/2$.
For all $p\ge13$ the
largest normalized
coordinate box occurs at
$p=13$, where the carrier
has $|m|<2$, $0<a<2$ and
$\lambda_pR_p>(3/20)2^p$.
With $t=1024$,
\[
 |m|+\tfrac a2(t+t^{-1})
 <2+1024+1/1024
 <(3/20)2^{13}
 \le\lambda_pR_p.
\]
The full ellipse, not a
sampled boundary point,
thus embeds for every
$p\ge13$, and
\[
 \pi\W(I_{\rm out})\le
 \frac{2\pi^2}{\log 1024}
 =\frac{\pi^2}{5\log2}<3
\]
using $\pi^2<10$
and $\log2>2/3$.
For $p=11,12$ the
simple exact sufficient
test $e^{7/2}<(33/20)^7
<169/5<17+12\sqrt2$
makes Lemma~\ref{lem:peripheral}
apply.  For the tail,
$e^3<27<17+12\sqrt2$
does the same.  All
four outer-incident
unoriented endpoint
types therefore vanish,
regardless of winding.

There are also disjoint
whole ellipses around
the individual $K_0,K_1,K_2$,
with $t$ respectively
$11,49,25$.  The exact
finite and all-$p$
physical-domain checks
include exclusion of
\emph{both other holes}
from each collar.
Equation \eqref{eq:collar}
gives the strict widths
\[
 \W(I_{K_0})<2.621,\quad
 \W(I_{K_1})<1.615,\quad
 \W(I_{K_2})<1.953.
\]
Under (H4), KKL Appendix A.2.3--A.2.5
bounds the total canonical
buffered weight, counting
returning arcs twice in
the local sums, by half
the sum of the peripheral
widths.  Since the outer
weights vanish,
\[
 0\le\lVert W_{\rm can}\rVert
 <\tfrac12(2.621+1.615+1.953)
 =3.0945<619/200.
\]
The remaining six
inner endpoint \emph{types}
can contain infinitely
many separate classes.
For any \emph{fixed} $p$,
KKL Appendix A.2.2 does imply
the nonconstructive support bound
$\#\{\gamma:q_\gamma>2\}
\le 3|\chi(Y_{K,p})|+4=10$;
this does not enumerate or give
a deck-word cutoff for any of them.
Thus
their classwise $q_\gamma$
and even the positivity
of the complete WAD
remain open.
\end{proof}

\section{Crest transport, nesting and global QC: precise limits}
\label{sec:limits}
\subsection{Finite collared four-push crest descendants}
The finite crest-transport problem uses
\emph{chosen nonmaximal circular parents},
not $P_p$ or $W_{\rm can}$.
At $p=3$ the parent is the native
Wave centered at $-9/64$,
$3/256<r<1/64$; for $p\ge4$
its chosen center is $-3x_p/4$
and radii $x_p/16<r<(p-1)x_p/(8p)$.
The parent conformal width is
$\log R_p^{\rm chosen}/(2\pi)$,
where $R_3^{\rm chosen}=4/3$
and $R_p^{\rm chosen}=2(p-1)/p$
for $p\ge4$.  The four
successive \emph{proper}
$\iota^*$ restrictions and
separate covering pushes
select two source terminal
families with endpoints
$(S_0,S_3)$ and
$(S_{17},S_0)$; their
target bases are
$(I_0,I_1,I_2,I_3)$
and $(I_7,I_8,I_9,I_0)$,
and their closed crests
$\overline{I_1\cup I_2}$
and $\overline{I_8\cup I_9}$
meet \emph{opposite sides}
$I_1,I_9$ of the \emph{same}
old edge $E_0$.
There is no assertion for
every third-stage path.

Let $\theta_p(r)=\arccos u_p(r)$,
where $u_p(r)\in(-1/8,1/4)$ is
the unique zero of
\[
 \Re G_p\!\left(d_p+r
 (u+i\sqrt{1-u^2})\right)=0.
\]
On a certified terminal
radius interval $(a,b)$
the conformal foliation
calculation gives the
\emph{selected} descendant
width and ratio
\begin{equation}\label{eq:crest}
 \W(S^\pm_{4,p})=
 \int_a^b\frac{dr}{r\,\theta_p(r)},
 \qquad
 \frac{\W(S^\pm_{4,p})}{\W(S_p)}
 =\frac{2\pi}{\log R_p^{\rm chosen}}
       \int_a^b\frac{dr}{r\,\theta_p(r)}.
\end{equation}
This follows from the
conformal logarithmic
strip with variable
angular height:
Cauchy--Schwarz
provides the extremal
density $1/\theta_p(e^x)$.
The archival certificate
bounds \emph{whole}
rational radius bins,
not sampled radii, using
the square-root-free
real/imaginary recurrence
\[
 X_{j+1}=X_j^2-Y^2B_j^2+c_p,\quad
 B_{j+1}=2X_jB_j,\quad
 Y^2=r^2(1-u^2).
\]
The relative cut and
full complex gate checks
are exact numerical
conditions; identification
with the selected KKL
proper-path subfamily
also requires (H2)
and the conformal
admissibility portion
of (H3).

For the maximal connected
\emph{circular collar only}
among those chosen
parent radii, the strict
finite outward enclosures
from the archive are:
\begin{center}
\begin{tabular}{rcc}
\toprule
$p$& lower bound $c_p^*$& upper bound $c_p^*$\\\midrule
3&0.98017&0.98038\\
4&0.62810&0.62818\\
5&0.52758&0.52764\\
6&0.47562&0.47566\\
8&0.43006&0.43010\\
12&0.39680&0.39684\\
13&0.39229&0.39233\\
16&0.38261&0.38265\\
\bottomrule
\end{tabular}
\end{center}
At $p=3$ the much
\emph{smaller fixed subcollar}
in the first transport
audit gives a \emph{different}
ratio in $(0.5676,0.5684)$;
it must not be conflated
with the full circular
optimum $(0.98017,0.98038)$.
The reported finite
interval enclosures concern
specifically selected collars;
their interpretation as KKL
proper-path widths is
conditional on (H2)--(H3),
not a canonical supremum
over all KKL Waves.

\begin{question}[Crest-transport limit]\label{q:crest}
Numerics suggest that
$c_p^*$ tends to the
positive constant
$c_\infty\approx0.3470049860435258$.
The proposed normalized
quadratic map is
$Q_\kappa(t)=1-\kappa t^2$,
where $\kappa\approx1.7903274919993457$
solves
$Q_\kappa^6(0)+Q_\kappa^3(0)=0$
on the distinguished branch.
The proposed first-order
law is
\[
 c_p^*=
 \frac{A_\infty}{\log(2-2/p)}
       +O(4^{-p}),\quad
 A_\infty\approx0.2405255277163131.
\]
Neither the convergence,
the $O(4^{-p})$ error,
nor a uniform all-period
selected q-side collar
survival is proved by
the finite audits.
The parameter asymptotic
\eqref{eq:tip-rate}
does \emph{not} alone
prove this stronger
four-stage
complex-domain and
prime-end statement.
\end{question}

\subsection{Depth two and the unique certified depth-three row}
For a \emph{realized} primitive
tuning word
$(p_1,\ldots,p_d)$
the depth-$j$ return is a
\emph{central degree-two
polynomial-like restriction}
of $f_c^{n_j}$ with
$n_j=\prod_{i=1}^j p_i$,
not the unrestricted
degree-$2^{n_j}$ polynomial.
For the nine \emph{candidate} depth-two
words $(p_1,p_2)\in\{3,4,5\}^2$,
the archived finite program reports
locally isolated real parameter
root cells, signed cut orders,
central degree-two returns and
a separately normalized third-return
Rouch\'e gap.  Their interpretation
as the \emph{named tuning suffixes}
requires a complete kneading and
straightening identification that
the following numbers alone do not
establish.  Subject also to (H1)--(H3),
the displayed numbers are proposed
strict lower witnesses and full-path
upper envelopes:
\begin{center}
\begin{tabular}{rcc}
\toprule
$(p_1,p_2)$ & $W^{\rm rad}$ lies in&
$Q_{\rm plane}$ lies in\\\midrule
$(3,3)$&(0.25224,0.25225)&(0.76475,0.76476)\\
$(3,4)$&(0.25854,0.25855)&(0.77250,0.77251)\\
$(3,5)$&(0.25999,0.26000)&(0.77430,0.77431)\\
$(4,3)$&(0.25179,0.25180)&(0.76418,0.76419)\\
$(4,4)$&(0.25848,0.25849)&(0.77243,0.77244)\\
$(4,5)$&(0.25999,0.26000)&(0.77431,0.77432)\\
$(5,3)$&(0.25170,0.25171)&(0.76406,0.76407)\\
$(5,4)$&(0.25847,0.25848)&(0.77242,0.77243)\\
$(5,5)$&(0.26000,0.26001)&(0.77431,0.77432)\\
\bottomrule
\end{tabular}
\end{center}
Consequently the inequality
$0.25<W^*<0.775$ is \emph{conditional}
on the geometric input and correct
suffix identification at these
nine physical depth-two charts.
The de Melo--van Strien rigidity
statement \cite[Theorem II.10.1]{deMeloVanStrien1993}
concerns quadratic maps with
equal \emph{complete} kneading
invariants.  It does \emph{not}
supply, from finite signed cut
tables alone, the hybrid
straightening or identification
of the chosen suffix.
Equating two independent
tip equations
$P_{p_1}=P_{p_1p_2}=0$
would be incorrect.
The data are finite exact
rational \emph{witness and
envelope} bounds, not
a characterization of
the optimum.

For depth three only the
\emph{candidate} word $(3,3,3)$
has the archived
complete finite numerical/interval
certificate; its realized
suffix interpretation still
requires the complete kneading
bridge:
$n_j=3,9,27$,
the nested target radii
$1/4,1/40,1/250$,
the three initial
Rouch\'e gaps strictly
$>0.127470,\ 0.006346,\
0.003900$, a \emph{separately
normalized} $81$-return
gap $>0.172673$,
and
\[
 W^{\rm rad}>0.252105,\qquad
 Q_{\rm plane}<0.764704.
\]
Consequently, subject to the
analogous geometric bridges
and identification of the nested
return, the physical marked
supremum would satisfy
$0.25<W^*<0.775$ \emph{for this one
depth-three row}.
None of these exact
finite recurrences proves
an all-word or infinite
tower theorem.

\subsection{Complex bounds, local hybrid maps and the missing global mark}
Lyubich--Yampolsky
\cite[p.~1220, Complex Bounds Theorem]{LyubichYampolsky1997}
gives an absolute positive
\emph{optimal} fundamental
annulus modulus for
every actually realized
finite non-doubling
renormalization of a real
quadratic, subject to its
explicit last-doubling
near-cusp exception.
This is \emph{not}
Theorem 1.1 restricted
to bounded essential period,
and it is not a modulus
bound for the specific
physical round gates
used here.
Douady--Hubbard
\cite[Theorem 1]{DH1985}
straightens each
appropriate quadratic-like
restriction \emph{locally}
around its little filled
Julia set.  Lyubich's draft
\cite[Lemma 4.33]{LyubichNotes}
states that a quadratic-like
map of fundamental modulus
at least $\delta>0$ admits
a $K(\delta)$-quasiconformal
hybrid straightening.
This statement and its
number were checked against
the actual draft (p.~143).
Neither a uniform modulus
for the selected physical
return gates nor a global
prime-end-marked QC map
follows from that lemma.  Neither
holomorphic-motion
extension
\cite{BersRoyden1986,Slodkowski1991}
nor Beltrami pullback
manufactures a missing
external prime-end marking
or an identification of
\emph{every} KKL A.1
and A.3.2 Wave class.

Here is an exact
conditional implication,
which makes the missing
hypothesis testable.
Let $\mathcal A_j$ be
the entire admissible
parent Wave class
on a physical depth-$j$
marked slit and
$\mathcal A_{\rm ref}$
a specified single-level
reference class.  Assume
for some \emph{common}
$C\ge1$:
\begin{itemize}
\item for every $\Gamma\in\mathcal A_j$
there is $\gamma\in\mathcal A_{\rm ref}$
with $\W(\Gamma)\le C\W(\gamma)$;
\item for every $\gamma\in\mathcal A_{\rm ref}$
there is $\Gamma\in\mathcal A_j$
with $\W(\gamma)\le C\W(\Gamma)$.
\end{itemize}
Taking suprema \emph{separately}
proves
\begin{equation}\label{eq:depth}
 \frac{W_{\rm ref}^*}{C}
 \le W_j^*\le C W_{\rm ref}^*.
\end{equation}
This \emph{two-sided
cofinal domination}
is strictly weaker than
the existence of a
global QC map.  A
globally defined real
equivariant $K'$-QC
marked conjugacy that
extends to the \emph{entire}
physical prime-end
completions, respects
$G_*$ and the set-valued
$\iota^*$, and
bijects all A.1/A.3.2
Wave classes would
imply it with $C=K'$
by extremal-length
quasi-invariance
\cite{AhlforsQC1966}.
\emph{Such a map has
not been constructed.}
If one \emph{assumes}
it with a depth-independent
$K'$, a general
single-level reference
gives
\[
 0.25/K'<W_j^*<0.78K',
\]
and a separately certified
reference with period
$p_{\rm ref}\ge13$ gives
$0.26/K'<W_j^*<0.78K'$.
These are \emph{conditional}
at every depth, with no
$K'^j$ factor because
the comparison is directly
to the reference.
For a positive \emph{matched}
parent and collared
four-push descendant,
their ratio is distorted
by at most $K'^2$
in either direction;
this does not apply
to two unrelated
optimizing suprema.
At infinite depth the
analogous liminf/limsup
bounds are \emph{closed},
and no limiting marked
slit or Wave convergence
is constructed.

A concrete obstruction
to \emph{deducing}
global marked $K'$
from local straightening
is furnished by
\[
 A_1=\{1<|z|<e^{2\pi}\},\qquad
 A_N=\{1<|z|<e^{2\pi N}\}.
\]
Both have an identical
real-symmetric conformal
collar near the inner
boundary, but their
full marked annuli
have moduli $1$ and
$N$.  Any QC map of
\emph{both marked
boundary components}
has $K\ge N$ by
modulus distortion.
This is a countermodel
to the \emph{logical
inference from local
data}, not a
counterexample in the
actual primitive-tip
KKL family.
The all-depth global
marked-QC transport
and an effective
$K'(\mu)$ are open.

\section{Limitations and disclosure}
The unconditional local
arithmetic/analytic claims
concern the \emph{selected
real} quadratic family.
Conclusions about the exact
physical KKL marked slit
surfaces and their complete
Wave families are
\emph{conditional} on the
geometric bridges of
Definition~\ref{def:bridges}.  The
local parameter
uniqueness does not
extend to all real
solutions or all
parameter components.
Neither the \emph{conditional}
positive unbuffered
independent Wave defect
nor the finite edge matrix
gives a positive
buffered WAD coefficient
or KKL \S5--6
amplification.  The
four-boundary
\emph{inner} buffered
support, maximizing
Waves, convergence of
canonical unrestricted
widths, the
four-push crest limit,
an all-word or
infinite-depth optimal
Wave theorem, and
uniform global
marked quasiconformal
transport remain open.
The family is real
and MLC in this
setting was already
settled by classical
theory; \emph{no
new MLC result} is
proved here.
The computation is
not a replacement
for independent
review of the
cited analytic proofs
or of this manuscript.
This text was
drafted and assembled
using AI assistance
from merged public
audit sources; it is
not yet independently
reviewed or
submitted as a
peer-reviewed
publication.

\section{Exact reproducibility, provenance and audit method}
\label{sec:reproduce}
The source of the finite statements in this
paper is the merged version of the repository
\url{https://github.com/dovvnloading/loops-research-archive}
as of 16 November 2026.
All source files cited here are \emph{read-only}
dependencies; this preprint lives entirely
under \texttt{papers/tip\_family\_preprint\_v2\_20261116/}
with one new dedicated workflow.
In particular no existing research
certificate, Lean proof, native-wave
source lock, result JSON or prior
CI workflow is modified.

\subsection{Exact verifiers and what they do}
Every key in the following displayed set
means the directory
\texttt{research/independent\_audits/KEY\_20261116/}.
For each of these directories the
new CI workflow runs
\texttt{python -B DIR/verify.py}
and discovers its adversarial
\texttt{test\_*.py} unit tests.
The whole theorem-chain verifier
\texttt{research/tip\_family\_theorem/verify\_chain.py}
and its unit tests are run separately.
\begingroup\fontsize{7}{8.5}\selectfont
\begin{verbatim}
real_primitive_tip_family_uniform
real_primitive_tip_family_unconditional
real_primitive_tip_family_all_p
real_primitive_tip_family_prime_end_identification
real_primitive_tip_family_canonical_width
real_primitive_tip_family_canonical_width_sandwich
tip_family_canonical_wave_gain
tip_family_limiting_wave_gain
tip_family_finite_p_wave_gain
tip_family_hypothesis_M_discharge
tip_family_buffered_wad
tip_family_renormalization_depth
tip_family_all_depths
tip_family_depth_uniform_qc
tip_family_global_marked_qc
crest_transport_tip_family
crest_transport_asymptotics
\end{verbatim}
\endgroup
The source-locked native-wave
reproductions are
\texttt{research/native\_wave\_transfer/regenerate.py}
and
\texttt{research/native\_wave\_transfer/regenerate\_period4.py},
both invoked with their
nonwriting \texttt{--check} option
where supported.  The
native \texttt{SPECIFICATION.md},
\texttt{RESULTS.md},
and experiment manifests
record the original
$p=3,4$ gate and
ideal marking.  They
are \emph{not} independent
mathematical proofs of
the KKL prime-end
identification.

The exact checks are divided
by mathematical function:
\begin{itemize}
\item The finite/unconditional
and all-period real auditors
check root-endpoint signs,
strict derivatives over
\emph{whole} parameter
intervals, signed
$6p$ orbit orders,
fixed-point brackets,
central return components,
full-complex-disk
Rouch\'e tail margins,
the $\kappa$-chart
rational inequalities,
and model-root isolation
(including six thousand
nonroot rational cells).
The prime-end auditor
checks the full
$3p$ disk gate
and discrete
$10/20$ labels
but \emph{not}
Carath\'eodory's theorem.
\item The width-sandwich,
canonical/limiting/finite
gain and M-discharge
auditors check entire
embedded parent
quadrilaterals,
source upper
cross-ratios,
the finite
missing-energy boxes,
the shrinking
normalized model
interval box
for all $p\ge13$,
the exact ordering
of physical source
versus target cuts,
and strict
reciprocal inequalities.
They do \emph{not}
synthesize or formally
verify a general
extremal-length proof.
\item The buffered auditor
checks \emph{whole
Joukowski ellipses},
disjointness from
each foreign hole,
peripheral-length
thresholds,
the rational
Landen comparison
and a strict
upper WAD norm.
It does not
enumerate the
inner ideal
arc classes.
\item The crest auditors
check entire
radius bins,
endpoint gates,
monotonicity,
imaginary-axis
crossings, the
source attachment
side, finite
integral enclosures
and the selected
parent denominator.
The depth auditors
check only the
specified finite
depth-two and
depth-three
realized charts,
and the QC
auditors check
the \emph{logical
conditional}
statements, not
a new global
QC construction.
\end{itemize}

\subsection{Why arithmetic is rigorous}
The main finite
primitive-root
packets have
\emph{adjacent}
dyadic endpoints
of denominator
$2^{256}$, obtained
from explicit
rational seeds and
sign bisection.
Every polynomial
interval evaluation
and derivative
uses a fixed
integer grid
$S=2^{512}$:
a rational lower
value is rounded
by \emph{floor}
and an upper
value by
\emph{ceiling}.
For interval
addition and
multiplication,
the four
endpoint products
are enclosed
with the same
outward rule;
interval squaring
explicitly uses
zero whenever
the input
straddles zero.
Inclusion
follows inductively
under each
polynomial
operation.
This works
equally for a
full complex
disk when
combined with
the analytic
majorant
recurrence
\eqref{eq:tail}.
The third-return
finite tail
uses a separately
outward $2^{-192}$
upper grid.
The original
crest checker
uses a $384$-bit
outward rational
interval grid
and $512$
\emph{whole}
radius bins,
with independent
bisection and
strict angular
monotonicity,
not floating
point quadrature.

For transcendental
bounds, Machin's
identity
$\pi=16\arctan(1/5)
-4\arctan(1/239)$
and rational
alternating arctangent
remainders give
certified $\pi$
brackets.
For positive
$x$ in the
specified domains
the atanh series
\[
 \log x=2\sum_{j=0}^{N-1}
 \frac{u^{2j+1}}{2j+1}
 +\mathcal{R}_N,\qquad
 u=\frac{x-1}{x+1},\quad
 |\mathcal{R}_N|\le
 \frac{2|u|^{2N+1}}
 {(2N+1)(1-u^2)}
\]
and alternating
trigonometric
Taylor bounds
give rational
enclosures for
logs, trigonometric
and inverse
trigonometric
functions.
The elliptic
$K(k)$ ratios
are enclosed by
outward rational
AGM or equivalent
power-series
bounds, with
explicit tails.
Every numerical
entry in a
strict theorem
is thus a
rational
inequality;
any human-readable
decimal endpoint
is a \emph{literal
rational number},
not a machine
float or a
measured approximation.
The all-$p$
tail steps
prove the
universal
monotonicity
or nested-box
property and
check its
worst $p=13$
or $p=17$
case by
rational
arithmetic.
They are
\emph{not}
finite sampling
extrapolations.

The proof
and code
trust boundary
is important.
KKL definitions
and Lemma A.11
\cite{KKL2026};
Douady--Hubbard's
straightening
\cite{DH1985};
McMullen's
almost-disjointness
\cite{McMullen1994};
the Misiurewicz
dendrite theorem
\cite{LyubichNotes};
Carath\'eodory's
extension
\cite{Caratheodory1913};
SC and
Riemann--Hurwitz
\cite{AhlforsInvariants,Forster1981}
are imported
\emph{by theorem
and explicit
hypotheses},
not certified
by Python or
the existing
finite Lean
incidence
models.
No novel
classical
theorem is
secretly
encoded as
a numeric
test.

\subsection{Build and CI contract}
From the
repository root:
\begin{verbatim}
bash papers/tip_family_preprint_v2_20261116/build.sh
python -B research/tip_family_theorem/verify_chain.py
python -B -m unittest discover -s \
  research/tip_family_theorem -p 'test_*.py' -v
\end{verbatim}
The dedicated
\texttt{tip-family-preprint-v2-20261116.yml}
GitHub Actions workflow also runs
the newer \texttt{tip\_family\_buffered\_wad\_inner}
audit, whose support-cardinality result
does \emph{not} enumerate the inner classes.
The workflow uses
plain
\texttt{pdflatex}
and \texttt{bibtex}
to produce and
upload the
\texttt{main.pdf}
artifact, and
rechecks all
seventeen named
verifiers and
their unit
tests.  Its
last job
dispatches the
unchanged
\texttt{lean-proofs.yml}
(Research validation)
and
\texttt{weighted-lean.yml}
(Weighted spectral
mathlib proofs)
\emph{at the
exact same
PR head SHA}
and demands
success from
both.  A
successful
build is not
an independent
mathematical
peer review,
nor proof of
the open
claims above.

\bibliographystyle{plain}
\bibliography{references}
\end{document}
